\documentclass[10pt,a4paper]{amsart}
\usepackage[margin=19mm]{geometry}
\usepackage{amsmath,amssymb,mathtools}
\usepackage[scr=rsfs,cal=euler]{mathalfa}
\usepackage{xfrac}
\usepackage[unicode,hidelinks,bookmarks=false]{hyperref}
\newcommand{\ie}{i.e.\ }
\newcommand{\cf}{cf.\ }
\newcommand{\resp}{respectively\ }
\newcommand{\Subsection}{\subsection}
\providecommand{\nobreakdash}{\nobreak\hbox{-}}
\setlength{\emergencystretch}{2em}
\multlinegap0pt
\usepackage{amssymb}
\usepackage{tikz}
\usepackage{xcolor}
\newtheorem{lemma}{Lemma}
\newcommand{\Gmn}{G_{m,n}}
\title{One construction for the Miura-ori flip-graph degree sequence}
\author{Chakshu Gupta}
\date{Source: arXiv:2607.05567v3 (CC BY 4.0)}
\begin{document}
\maketitle

\noindent\emph{Source and licence.} One construction for the Miura-ori flip-graph degree sequence. Version arXiv:2607.05567v3 (CC BY 4.0), distributed by arXiv under the Creative Commons Attribution 4.0 International licence, http://creativecommons.org/licenses/by/4.0/, as recorded in the arXiv licence metadata for this exact version and shown on the abstract page https://arxiv.org/abs/2607.05567v3. Full licence notice: \texttt{licenses/arxiv-2607.05567-CC-BY-4.0.txt}.\par\smallskip

\noindent\textit{Editorial scope.} Counted unit: the article's lemma lem:frozenbdy (source0.tex lines 1463--1465). The article's complete original proof of it is delivered with it (source0.tex lines 1467--1529, one \texttt{proof} environment of 63 lines). No mathematics is written, repaired or re-derived: the statement and the proof are the article's own lines, in the article's own notation, and no retained line is substituted or re-flowed.  Retained uncounted as the source-local context the proof needs: lines 1363--1366.  Nothing was omitted from any retained span, so the package carries no omission record.  No retained line contains a citation, so the wrapper carries no bibliography and no bibliography entry.  No retained line contains a figure environment, an image-inclusion command or drawing code, so no image is delivered and the package declares no asset dependency.  Editorial formatting only, in addition: an AMS \texttt{amsart} preamble that restates the article's own declarations and macros used by the retained lines, namely the environments \texttt{lemma} on separate counters, together with the standard packages those lines need.  The retained environments print in this document as Lemma 1 in order of appearance, whereas the article prints the same items with the numbering of its own full document.  The complete native source of the article as distributed in the arXiv e-print for this exact version is bundled unchanged as \texttt{sources/204554/source0.tex}.
\begin{lemma}[Boundary extremum]\label{lem:boundary}
In every height function on $\Gmn$ with $m,n\ge2$, the first and last columns each
carry a strict local extremum.
\end{lemma}

\begin{lemma}[Frozen single-extremum boundary]\label{lem:frozenbdy}
On the count block, $\mathbf u_1^{\!\top}Q=0$ and $Q\mathbf v_1=0$.
\end{lemma}

\begin{proof}
By left--right symmetry, it suffices to prove $\mathbf u_1^{\!\top}T_0=\mathbf u_1^{\!\top}$ on the
count block. Its mirror identity $T_0\mathbf v_1=\mathbf v_1$
then holds, and $\mathbf u_1^{\!\top}Q=0$ and $Q\mathbf v_1=0$ follow. Written on
the states, $\mathbf u_1^{\!\top}T_0=\mathbf u_1^{\!\top}$ concerns the columns $a$ that may sit to
the left of $b$. A column $a$ is \emph{valid} when it is admissible with $b$, carries one
extremum with right neighbour $b$, and leaves $b$ extremum-free in the triple
$(a,b,c)$. The claim is that for every admissible pair $(b,c)$,
\begin{equation}\label{eq:frozenid}
  \#\{\text{valid }a\}=[\,b\ \text{has one extremum with right neighbour }c\,].
\end{equation}
A left neighbour can only destroy extrema of $b$, since an extremum needs its
present neighbours to share a colour and a new neighbour may break that. So a valid
$a$ takes the third colour at each extremum cell of $b$, lying above $b$ at a
maximum and below at a minimum.

\medskip\noindent
The columns $a$ and $b$ lift to height walks $\alpha,\beta$ as in the proof of
Lemma~\ref{lem:boundary}, with steps $\varepsilon_i:=\alpha_{i+1}-\alpha_i$ and
offset $\eta_i:=\beta_i-\alpha_i$, both in $\{\pm1\}$, and $\beta$ again a $\pm1$
walk. That proof's flip rule lets $\eta$ change from $-1$ to $+1$
only on a descent, and from $+1$ to $-1$ only on an ascent.

\medskip\noindent\emph{Every valid $a$ is a uniform shift.} A valid $a$ has one
extremum; say a maximum at $p$, so $\eta_p=-1$. Then $\alpha$ has a single local
maximum. Suppose instead a second maximum $p''>p$, with a local minimum $\mu$
between them; the case $p''<p$ is symmetric. On the descent $[p,\mu]$, the flip
rule bars every $\eta$-change $+1\to-1$. So either $\eta$ reaches $+1$ before $\mu$,
making the walk minimum $\mu$ an extremum of $a$, as $\eta_\mu=+1$ there; or $\eta$
stays $-1$ through $\mu$, and on the ascent $[\mu,p'']$ the flip rule bars
$-1\to+1$, so $\eta_{p''}=-1$ makes $p''$ an extremum. Either way, $a$ has a second
extremum, against validity. So $\alpha$ is unimodal, with $\varepsilon_i=+1$ for
$i<p$ and $-1$ for $i\ge p$. Then the increment
\[
  \eta_{i+1}-\eta_i=(\beta_{i+1}-\beta_i)-\varepsilon_i
\]
is $\le0$ for $i<p$ and $\ge0$ for $i\ge p$, so $\eta$ falls to $\eta_p=-1$ at the peak and rises after. It
remains to pin both corners at $-1$. Consider cell $0$. If $p=0$, then
$\eta_0=\eta_p=-1$. Otherwise, $p>0$ gives $\varepsilon_0=+1$ by unimodality, and
cell $0$ is not an extremum, so the test $\varepsilon_0=\eta_0$ of
Lemma~\ref{lem:boundary} fails, forcing $\eta_0=-1$. Cell $m-1$ is the same. If
$p=m-1$, then $\eta_{m-1}=-1$. Otherwise, $\varepsilon_{m-2}=-1$ and cell $m-1$ is
not an extremum, so $\varepsilon_{m-2}=-\eta_{m-1}$ fails, forcing $\eta_{m-1}=-1$.
A $\pm1$ sequence non-increasing to $-1$ then non-decreasing back
to $-1$ is constant, so $\eta\equiv-1$ and $a=b+1$. A minimum gives $a=b-1$.

\medskip\noindent\emph{The shift count.} The column $b+1$ lies above $b$, so it
destroys exactly the maxima of $b$ and leaves its minima. Its own extrema against
$b$ are the local maxima of $\beta$, endpoints included and independent of $c$;
likewise, $b-1$ destroys the minima and has extrema at the local minima. A shift is
valid only if it destroys every extremum of $b$, so $b+1$ can be valid only when
$b$ has no minimum, and $b-1$ only when $b$ has no maximum. By
Lemma~\ref{lem:boundary}, $b$ carries at least one extremum. Suppose it carries
exactly one, a maximum. Then $b$ has no minimum, and the same unimodality applied
to $b$ against $c$ gives $\beta$ a single local maximum; so $b+1$ has one extremum
and is valid, while $b-1$ leaves that maximum and is not. A single minimum is
symmetric. Suppose instead $b$ carries two or more extrema. Then either it has both
a maximum and a minimum, so neither shift destroys all of them; or it has two of
one type, say two maxima, so $b+1$ inherits two extrema from the two local maxima
of $\beta$. Either way, no shift is at once valid and single-extremum. The valid
columns therefore number $[\,b\ \text{has one extremum with right neighbour }c\,]$,
which is~\eqref{eq:frozenid}.
\end{proof}

\end{document}
